\chapter{La mitad que falta}
\label{cap:mitad-que-falta}

El capítulo anterior demostró que $G$ no tiene demostración. Para la indecidibilidad falta la otra
mitad: que tampoco tiene refutación. Este capítulo la construye hasta donde llega —que es lejos— y
después hace algo que un libro normal evitaría: explicar con precisión qué es lo que falta, por qué
no se puede coger el atajo, y cuánto trabajo hay hecho.

Porque el interés de esta mitad no es que esté incompleta. Es \emph{dónde} se atasca.

\section{El argumento, hasta donde llega}

Supóngase que $\lnot G$ sí fuera demostrable.

\begin{matematica}[Los cuatro pasos, y el quinto]
\begin{enumerate}
\item De $\Prf(\lnot G)$ se pasa a $\text{\ident{axioms}} \derives \lnot G$ por el puente.
\item El punto fijo, leído de derecha a izquierda, dice
  $\lnot\Prov(\codigo{G}) \Rightarrow G$. Con (1), la teoría demuestra
  $\lnot\lnot\Prov(\codigo{G})$.
\item Por eliminación de la doble negación, la teoría demuestra $\Prov(\codigo{G})$.
\item Y ahora hace falta el paso que no es lógica: de que \emph{la teoría demuestre que $G$ es
  demostrable} concluir que $G$ \textbf{realmente} lo es. Con eso, contradicción con el capítulo
  anterior.
\end{enumerate}
\end{matematica}

El paso (4) es todo el problema. Los tres primeros son manipulación formal; el cuarto dice que la
teoría no se equivoca al afirmar que algo es demostrable. Y eso no se sigue de que sea consistente.

\section{La hipótesis, con nombre}

El proyecto no la esconde: la declara.

\leanfrag{Reflects}

\selee{«\ident{Reflects} $G$ dice: si la teoría demuestra que $G$ es demostrable, entonces $G$ tiene
demostración»}

Es la \textbf{representabilidad negativa}, el espejo de D1. D1 dice que si algo es demostrable, la
teoría lo sabe; la reflexión dice que si la teoría lo afirma, es verdad. La primera es un teorema
del capítulo~\ref{cap:d1}. La segunda no lo es, y no puede serlo.

Con ella delante, la mitad se cierra:

\leanfrag{goedel_first_unrefutable}

\leanfrag{goedel_first_undecidable}

\selee{«bajo consistencia y reflexión, $G$ no es demostrable ni refutable»}

\section{Por qué el atajo está cerrado — y quien lo cierra es Gödel}

La tentación evidente es demostrar la reflexión dentro de la teoría. Haría falta que, para cada
$\varphi$ indemostrable, la teoría probara $\lnot\Prov(\codigo{\varphi})$. Los dos casos que uno
probaría primero son los dos que lo cierran:

\begin{muro}[Los dos callejones, y son el mismo]
Tómese $\varphi = \bot$. Si la teoría es consistente, $\bot$ es indemostrable, así que haría falta
$\derives \lnot\Prov(\codigo{\bot})$. Pero esa fórmula \textbf{es} $\mathrm{Con}(T)$, la sentencia
de consistencia — y por el segundo teorema de Gödel la teoría no la demuestra.

Tómese $\varphi = G$. Por el punto fijo, $\derives \lnot\Prov(\codigo{G})$ equivale a
$\derives G$ — y el capítulo anterior acaba de demostrar que eso no pasa.

La vía «que la teoría refute la demostrabilidad» no está cerrada por falta de trabajo:
\textbf{está cerrada por los teoremas de este libro}.
\end{muro}

\begin{leccion}[Un obstáculo que se demuestra es información, no un hueco]
Merece la pena parar en la forma de ese argumento. No dice «no lo hemos conseguido»: dice
\emph{dónde} está el obstáculo y \emph{qué} teorema lo pone. Un lector puede comprobarlo. Y como
efecto lateral, explica por qué existe Rosser: para tener las dos mitades desde consistencia simple
hay que cambiar de sentencia, porque con \emph{esta} sentencia el camino corto lo corta Gödel~II.
\end{leccion}

\section{El postulado que escondía todo esto}

Esta mitad estuvo dada por cerrada durante tres meses, y no lo estaba.

\begin{refutado}[Lo que decía la capa antigua]
El 13 de junio de 2026 la mitad $\noderives\lnot G$ se cerró apoyándose en un postulado escrito
como \textbf{bicondicional}:
\[ \bigl(\text{\ident{axioms}} \derives \Prov(\codigo{\varphi})\bigr)
   \;\Longleftrightarrow\;
   \bigl(\text{\ident{axioms}} \derives \varphi\bigr). \]
La dirección de derecha a izquierda es D1, y es un teorema. La de izquierda a derecha \textbf{es la
reflexión}, y es falsa en general. Y el teorema que se apoyaba en ella estaba enunciado bajo
\emph{consistencia simple}: es decir, afirmaba más de lo que Gödel permite. Se retiró.
\end{refutado}

\begin{leccion}[El mismo modo de fallo, por tercera vez]
Un enunciado plausible, aceptado como postulado, que resulta ser más fuerte de lo que se puede
sostener. Ya pasó con un axioma del kernel (capítulo~\ref{cap:kernel}) y con el axioma que hacía
inconsistente la teoría (capítulo~\ref{cap:numerales}). Aquí no produjo una contradicción: produjo
algo más difícil de detectar, \textbf{un teorema que decía de más}. El compilador no protesta
cuando una hipótesis es demasiado fuerte — sólo cuando falta.

La reparación no fue demostrar el postulado. Fue \textbf{partirlo}: quedarse con la mitad que es
teorema y dejar la otra a la vista como hipótesis.
\end{leccion}

\section{Cómo se descarga: dos piezas}

La reflexión no es indemostrable en absoluto — es indemostrable \emph{dentro} de la teoría. Desde
la metateoría se descarga, y el proyecto tiene el camino escrito:

\leanfrag{reflects_of_omega}

Hacen falta dos cosas. La primera es la \textbf{$\omega$-consistencia} —la de verdad, la que el
capítulo anterior distinguía del nombre parecido—. La segunda es un verificador que además de
aceptar las derivaciones buenas \textbf{rechace} las malas:

\leanfrag{NegVerifier}

\selee{«para toda $\varphi$ no demostrable y toda cadena estándar, la teoría refuta que esa cadena
verifique $\varphi$»}

Es el espejo negativo de D1: D1 dice que el verificador acepta lo que debe; esto dice que rechaza
lo que debe. Con las dos, la indecidibilidad sale sin hipótesis meta:

\leanfrag{goedel_first_undecidable_omega}

\begin{observacion}
La demostración de \ident{reflects_of_omega} es donde vive el \textbf{único uso esencial de lógica
clásica} de todo el proyecto, y el capítulo~\ref{cap:objeto-y-meta} lo identificó: es el principio
de Markov. Tiene sentido que esté aquí. Lo que se hace es concluir «existe una demostración» de la
imposibilidad de que no exista, y eso —sobre un enunciado $\Sigma_1$ con matriz decidible— es
exactamente Markov.
\end{observacion}

\section{Y ahora la cifra incómoda}

\ident{NegVerifier} está escrito. Está enunciado con precisión, tiene su papel claro en la cadena y
dos teoremas lo consumen como hipótesis.

\textbf{Y no había ni un teorema en producción que lo concluyera.} Así se midió al escribir este
capítulo, la mañana del 10 de septiembre de 2026: cero.

\begin{observacion}[Corregido esa misma tarde, y las dos mitades de la corrección importan]
La frase de arriba se quedó a medias en unas horas, y la parte que se cayó no es la que parece.

Lo que \textbf{sigue siendo cierto} es que ningún teorema concluye \ident{NegVerifier} de manera
incondicional. Lo que \textbf{dejó de serlo}: ya no es cierto que no haya nada. Existe
\ident{negVerifier_of_deudas}, que lo concluye \textbf{a partir de dos obligaciones enunciadas},
\ident{DEUDA_chainNeg} y \ident{DEUDA_inNeg}; y la solidez estructural del verificador —el módulo
que el plan llamaba «el corazón» y estimaba en 300--500 líneas— resultó estar \textbf{ya
demostrada}: \ident{verifier_sound} \emph{es} \ident{decodeChain_prf}.

Y lo que era \textbf{directamente falso} es la causa que esta sección daba: «lo que lo bloquea
—la inyectividad de la codificación numérica— vive en un sondeo, fuera del build». No: \ident{consN_inj}
y \ident{codeNat_ne} están en producción desde el 2 de septiembre. Lo dice la propia auditoría de
este libro, y es un buen ejemplo de por qué §2.6 vale también \emph{hacia dentro}: la frase se
copió de un plan del proyecto en vez de medirse.

El bloqueo real era otro, y hubo que medirlo: la clase de testigos admitía basura cuya refutación
exigía \textbf{evaluar el emparejamiento de Cantor}, que es inviable. Se estrechó
(\textsc{adr}-022), con el precio escrito —la $\omega$-consistencia queda algo más fuerte— y con
la garantía que lo hace admisible: el código de \emph{cualquier} derivación real sigue siendo un
testigo de la clase.
\end{observacion}

\avisomediagodel

\begin{leccion}[Qué es exactamente lo que falta, y por qué decirlo así vale más]
Compárese con la manera habitual de dejar esto: «la otra mitad es análoga» o, peor, «Gödel~I,
completo». Lo que este proyecto puede decir en su lugar es:

\begin{center}\small
\begin{tabularx}{\linewidth}{@{}l >{\raggedright\arraybackslash}X@{}}
\toprule
pieza & estado \\
\midrule
el argumento $\noderives\lnot G$ desde la reflexión & \textbf{demostrado} \\
la reducción reflexión $\leftarrow$ $\omega$-consistencia $+$ verificador negativo & \textbf{demostrada} \\
el verificador negativo & \textbf{reducido a dos obligaciones con nombre} \\
\bottomrule
\end{tabularx}
\end{center}

Eso no es un hueco: es un plano con una pieza marcada en rojo. La diferencia entre las dos maneras
de escribirlo es la diferencia entre un libro que se puede auditar y uno que hay que creer.
\end{leccion}
